% ! TEX root = ../main.tex

\newlength{\tabOrcaYwidth}
\setlength{\tabOrcaYwidth}{\dimexpr(\textwidth - 1.0in - 1.5in - 8\tabcolsep) / 2\relax}

\newcommand{\tabOrcaRelatedInner}[1]{%
  \begin{tabular}
    [t]{@{}>{\raggedright\arraybackslash}p{1.5in}>{\raggedright\arraybackslash}p{\tabOrcaYwidth}>{\raggedright\arraybackslash}p{\tabOrcaYwidth}@{}}
    #1
  \end{tabular}%
}

\newcommand{\tabOrcaGreyhound}{%
  \tabOrcaRelatedInner{
    \textbullet~Iteration boundary detection & \textbullet~Autocorrelation over collectives & \textbullet~\texttt{PostTimestepAdvance} \\
    \textbullet~Fail-slow detection & \textbullet~Track iteration times & \textbullet~Track iteration times \\
    \textbullet~Straggler localization & \textbullet~Aggregate NCCL events (off-fabric) & \textbullet~\texttt{GROUPBY swid} (on-fabric) \\
    \textbullet~Straggler cause disambiguation & \textbullet~Microbenchmarks & \textbullet~Microbenchmarks or trace events \\
    \textbullet~Online interventions & \textbullet~Tune micro-batch distribution or parallelism & \textbullet~Supported via control plane \\
  }%
}

\newcommand{\tabOrcaFlare}{%
  \tabOrcaRelatedInner{
    \textbullet~Stall detection/disambiguation & \textbullet~Call stack analysis & \textbullet~In-situ reductions \\
    \textbullet~Slowdown analysis & \textbullet~Periodic metrics & \textbullet~In-situ trace event aggregation \\
  }%
}

\newcommand{\tabOrcaAegis}{%
  \tabOrcaRelatedInner{
    \textbullet~Fail-stop hardware & \textbullet~Per-node NCCL counters pulled on timeouts & \textbullet~Workload-level trace events \\
  }%
}

\newcommand{\tabOrcaHolmes}{%
  \tabOrcaRelatedInner{
    \textbullet~Outlier collective detection & \textbullet~Random forest model & \textbullet~Policy-dependent \\
    \textbullet~Localization & \textbullet~Graph search & \textbullet~Learn collective topology at bootstrap, \texttt{GROUPBY~swid}, trace critical path in reverse \\
  }%
}

\newcommand{\tabOrcaMycroft}{%
  \tabOrcaRelatedInner{
    \textbullet~Straggler detection/localization & \textbullet~Fine-grained trace NCCL events buffered locally, flushed to Kafka if heuristics trigger & \textbullet~Aligned dataframes analyzed in-situ \\
  }%
}

\newcommand{\tabOrcaTrainCheck}{%
  \tabOrcaRelatedInner{
    \textbullet~Correctness invariants & \textbullet~Offline trace analysis for invariant generation, online enforcement & \textbullet~Timestep-aligned traces for offline generation, in-situ enforcement via OrcaFlow \\
  }%
}

\newcommand{\tabOrcaReveal}{%
  \tabOrcaRelatedInner{
    \textbullet~Anomaly detection and attribution & \textbullet~Time-windowed metrics + unsupervised learning & \textbullet~Workload-level trace events \\
  }%
}

\newcommand{\tabOrcaMinder}{%
  \tabOrcaRelatedInner{
    \textbullet~Fail-slow hardware & \textbullet~Cluster-level metrics and anomaly detection & \textbullet~Workload-level trace events \\
  }%
}

\newcommand{\tabOrcaPerfTracker}{%
  \tabOrcaRelatedInner{
    \textbullet~Anomaly detection and attribution & \textbullet~Always-on iteration monitoring & \textbullet~Identical workflow, but workload-level \\
    & \textbullet~Trigger-driven detailed capture & \\
    & \textbullet~In-situ telemetry reductions & \\
  }%
}

\appendices
\section{ML Training Observability Comparison}\label{sec:appdx}

\newcommand{\tabOrcaRelated}{%
\begin{table*}[t]
  \hyphenpenalty=10000
  \exhyphenpenalty=10000
  \centering
  \scriptsize
  \newcolumntype{Y}{>{\raggedright\arraybackslash}X}
  \begin{tabularx}{\textwidth}{>{\raggedright\arraybackslash}p{1.0in} X}
    \toprule
    \textbf{System}                                                     &
    \tabOrcaRelatedInner{
    \textbf{Diagnostic Task}                                            & \textbf{Their Approach} & \textbf{ORCA Equivalent} \\
    }
    \\
    \midrule
    Greyhound~\cite{Wu2025:GREYHOUND} \emph{(Fail-slows)}               &
    \tabOrcaGreyhound                                                                                                        \\
    \midrule
    Flare~\cite{Cui2026:Flare} \emph{(Training~anomalies)}              &
    \tabOrcaFlare                                                                                                            \\
    \midrule
    Holmes~\cite{Yao2025:Holmes}\newline \emph{(Fail-slows)}                    &
    \tabOrcaHolmes                                                                                                           \\
    \midrule
    Mycroft~\cite{Deng2025:Mycroft} \emph{(Training~anomalies)}         &
    \tabOrcaMycroft                                                                                                          \\
    \midrule
    TrainCheck~\cite{Jiang2025:TrainCheck} \emph{(Training~correctness)} &
    \tabOrcaTrainCheck                                                                                                       \\
    \midrule
    Reveal~\cite{Chen2025:Reveal} \emph{(Training~anomalies)}           &
    \tabOrcaReveal                                                                                                           \\
    \midrule
    Minder\textsuperscript{\#}~\cite{Deng2025:Minder} \emph{(Fault~diagnosis)} &
    \tabOrcaMinder                                                                                                           \\
    \midrule
    Aegis\textsuperscript{\#}~\cite{Dong2025:Aegis} \emph{(Fault~diagnosis)}                &
    \tabOrcaAegis                                                                                                            \\
    \midrule
    PerfTracker\textsuperscript{\#}~\cite{Guan2025:PerfTracker} \emph{(Training~anomalies)} &
    \tabOrcaPerfTracker                                                                                                      \\
    \bottomrule
  \end{tabularx}
  {\par\vskip 0.2em\raggedright\scriptsize\itshape \textsuperscript{\#} Minder, Aegis, and PerfTracker target multi-tenant environments. ORCA currently supports a single large workload.\par}
  \caption{Diagnostic approaches in ML training observability and their ORCA equivalents.}\label{tab:orca:related-work}
\end{table*}%
}

\tabOrcaRelated{}

Recent ML training systems build application-specific pipelines spanning
straggler localization~\cite{Wu2025:GREYHOUND,Yao2025:Holmes,Deng2025:Mycroft},
fault diagnosis~\cite{Deng2025:Minder,Dong2025:Aegis}, and correctness
enforcement~\cite{Jiang2025:TrainCheck}.
ORCA generalizes the pattern.
Holmes~\cite{Yao2025:Holmes} uses graph search over collective telemetry to
find anomalies.
ORCA's aligned collection makes stragglers surface through
simple grouped aggregations, no complex detection needed.
Mycroft~\cite{Deng2025:Mycroft} buffers fine-grained NCCL~\cite{:NCCL} telemetry locally
and flushes to Kafka~\cite{Kreps2011:Kafka} on trigger, but notes the scalability limits of central
logging.
ORCA's tree-based overlay scales out naturally.
TrainCheck~\cite{Jiang2025:TrainCheck} enforces training correctness via
offline invariant generation and online checks.
ORCA's timestep-aligned columnar analytics accelerate both stages.
Minder~\cite{Deng2025:Minder},
Aegis~\cite{Dong2025:Aegis}, and PerfTracker~\cite{Guan2025:PerfTracker} operate at the cluster
infrastructure level,
using hardware counters, metrics, and fine-grained trace events to identify faulty nodes across
workloads.
Their scopes range from identifying bad hardware across many workloads treated as black boxes to
providing tuning recommendations to customers.
ORCA provides
high-fidelity visibility into a single workload but does not yet support
cross-workload diagnosis in multi-tenant environments.
\Cref{tab:orca:related-work} summarizes the diagnostic mechanisms implemented in these systems and
their ORCA equivalents.

Among these systems, Greyhound~\cite{Wu2025:GREYHOUND} and PerfTracker~\cite{Guan2025:PerfTracker}
are the closest to ORCA in scope.
Greyhound implements a complete feedback loop: iteration times are
tracked, NCCL events are collected for localization when outliers are
detected, microbenchmarks are dispatched to validate stragglers, and online
tuning rebalances work in response.
This is a fixed-function version of what ORCA expresses as composable operators.
ORCA's timestep dataframe abstraction preserves fine-grained trace events,
enabling disambiguation of stragglers without needing microbenchmarks for validation.
Greyhound's control semantics require pausing the application,
while ORCA implements an asynchronous control plane with atomicity and timestep-consistency
guarantees.
Additionally, ORCA's on-fabric broadcast tree enables faster message dissemination than a central
controller.
PerfTracker implements a two-level loop: lightweight always-on monitoring of iteration times, with
detailed capture triggered by outliers.
Node-level aggregation reduces 300\,MB of telemetry to 30\,KB of aggregated signatures.
Long symbol names are noted as contributing to size amplification, which ORCA avoids via
dictionary-encoded Arrow columns.